\documentclass[9pt,a4paper,twocolumn]{ltjsarticle}
\usepackage[top=20mm,bottom=22mm,left=17mm,right=17mm,columnsep=8mm]{geometry}
\usepackage{amsmath}
\usepackage{booktabs}
\usepackage[hidelinks]{hyperref}
\newcommand{\NSamples}{442}
\newcommand{\NFeatures}{10}
\newcommand{\NTrain}{353}
\newcommand{\NTest}{89}
\newcommand{\Seed}{0}
\newcommand{\TargetMean}{152.1}
\newcommand{\TargetStd}{77.0}
\newcommand{\TrainMse}{2734.8}
\newcommand{\TrainRmse}{52.3}
\newcommand{\TrainMae}{42.6}
\newcommand{\TrainRsq}{0.554}
\newcommand{\TestMse}{3424.3}
\newcommand{\TestRmse}{58.5}
\newcommand{\TestMae}{46.2}
\newcommand{\TestRsq}{0.332}
\newcommand{\CpuModel}{Intel Core i5-10400 2.90GHz}
\newcommand{\CpuCores}{12}
\newcommand{\MemoryGib}{7.7}
\newcommand{\PythonVersion}{3.13.3}
\newcommand{\NumpyVersion}{2.3.4}
\newcommand{\SklearnVersion}{1.7.2}

% 各文の ID。claims.json の sentences の ID と同じ
\newcommand{\sid}[1]{\textsuperscript{\textsf{\scriptsize[#1]}}}
\newcommand{\RepositoryUrl}{https://github.com/auto-res2/diabetes-linear-regression-6}
% 論文の途中からの抜粋: 節と頁の番号を途中から始める
\setcounter{section}{3}
\setcounter{page}{3}
\setcounter{equation}{2}
\setcounter{table}{1}
\setcounter{footnote}{1}

\begin{document}

\twocolumn[{
\begin{center}
{\Large\bfseries 糖尿病データセットに対する線形回帰モデルの予測性能}\\[2mm]
{\normalsize auto-res2}\\[1mm]
{\small （「4 実験設定」および「5 結果」の抜粋）}
\end{center}
\vspace{3mm}
}]

\section{実験設定}

\subsection{データセット}
\sid{R01}Efron ら~\cite{efron2004} が用いた糖尿病データセットを使用した．\sid{R02}本データセットは糖尿病患者 \NSamples{} 名の記録からなり，各患者について年齢，性別，BMI，平均血圧，および 6 種類の血清測定値の計 \NFeatures{} 個の説明変数と，ベースラインから 1 年後の病勢進行度を表す目的変数が与えられている．\sid{R03}データは公開されているファイルから取得し，ファイルの SHA-256 ハッシュ値が事前に固定した値と一致することを確認した上で使用した．\sid{R04}説明変数には標準化などの前処理を施していない．\sid{R05}目的変数の平均は \TargetMean{}，標準偏差は \TargetStd{} である．

\subsection{データの分割}
\sid{R06}全 \NSamples{} 件を，訓練データ \NTrain{} 件とテストデータ \NTest{} 件に無作為に分割した．\sid{R07}テストデータの割合は 20\% とし，乱数のシードは \Seed{} に固定した．

\subsection{モデル}
\sid{R08}説明変数を $\boldsymbol{x} = (x_1, \dots, x_p)^\top$（$p = \NFeatures$），目的変数を $y$ とし，線形回帰モデル
\begin{equation}
  \hat{y} = \beta_0 + \sum_{j=1}^{p} \beta_j x_j
  \label{eq:model}
\end{equation}
を，訓練データに対して最小二乗法により当てはめた．\sid{R09}すべての説明変数を用い，正則化は行っていない．\sid{R10}実装には scikit-learn~\cite{sklearn} の \texttt{LinearRegression} を用いた．

\subsection{評価指標}
\sid{R11}予測性能は，平均二乗誤差（MSE），二乗平均平方根誤差（RMSE），平均絶対誤差（MAE），および決定係数（$R^2$）により評価した．\sid{R12}$n$ 件のデータに対する MSE と $R^2$ は
\begin{align}
  \mathrm{MSE} &= \frac{1}{n} \sum_{i=1}^{n} (y_i - \hat{y}_i)^2, \\
  R^2 &= 1 - \frac{\sum_{i=1}^{n} (y_i - \hat{y}_i)^2}{\sum_{i=1}^{n} (y_i - \bar{y})^2}
\end{align}
で定義される．\sid{R13}ここで $\bar{y}$ は当該データにおける目的変数の平均である．\sid{R14}RMSE は MSE の平方根，MAE は $|y_i - \hat{y}_i|$ の平均である．\sid{R15}各指標は訓練データとテストデータのそれぞれについて算出した．

\subsection{計算環境}
\sid{R16}実験は \CpuModel{}（論理 \CpuCores{} コア），メモリ \MemoryGib{}\,GiB の計算機上で行い，GPU は使用していない．\sid{R17}ソフトウェアには Python \PythonVersion{}，NumPy \NumpyVersion{}，scikit-learn \SklearnVersion{} を用いた．\sid{R18}実験コードと出力は公開している\footnote{\url{\RepositoryUrl}}．

\section{結果}

\sid{R19}表~\ref{tab:metrics} に，訓練データおよびテストデータに対する各評価指標の値を示す．\sid{R20}テストデータに対する MSE は \TestMse{}，RMSE は \TestRmse{}，$R^2$ は \TestRsq{} であった．\sid{R21}テストデータに対する RMSE は目的変数の標準偏差 \TargetStd{} を下回っており，線形回帰モデルは目的変数の平均値を一律に予測する場合よりも小さい誤差で病勢進行度を予測できている．

\begin{table}[t]
\centering
\caption{\sid{R26}線形回帰モデルの評価指標．MSE，RMSE，MAE は小さいほど，$R^2$ は大きいほど良い．}
\label{tab:metrics}
\begin{tabular}{lrrrr}
\toprule
データ & MSE & RMSE & MAE & $R^2$ \\
\midrule
訓練（\NTrain{} 件） & \TrainMse & \TrainRmse & \TrainMae & \TrainRsq \\
テスト（\NTest{} 件） & \TestMse & \TestRmse & \TestMae & \TestRsq \\
\bottomrule
\end{tabular}
\end{table}

\sid{R22}一方，テストデータに対する $R^2$ は訓練データに対する値 \TrainRsq{} を大きく下回った．\sid{R23}MSE についても，テストデータでの値は訓練データでの値 \TrainMse{} より大きい．\sid{R24}これらの結果は，本モデルが訓練データに対してはある程度当てはまるものの，未知のデータに対する説明力は限定的であることを示している．\sid{R25}説明変数と目的変数の関係のうち，式~\eqref{eq:model} の線形モデルで捉えられるのは一部にとどまると考えられる．

\begin{thebibliography}{9}
\small
\bibitem{efron2004} B.~Efron, T.~Hastie, I.~Johnstone, and R.~Tibshirani: Least angle regression, \textit{The Annals of Statistics}, Vol.~32, No.~2, pp.~407--499 (2004).
\bibitem{sklearn} F.~Pedregosa et al.: Scikit-learn: Machine learning in Python, \textit{Journal of Machine Learning Research}, Vol.~12, pp.~2825--2830 (2011).
\end{thebibliography}

\end{document}
